\section{扰动响应预测}
% 4
\begin{frame}[t]{扰动响应预测：从简单基线到细胞基础模型}
\vspace{.05cm}
\centering
\begin{tikzpicture}[stp/.style={box,text width=2.55cm,inner sep=7pt,minimum height=1.15cm},
  nt/.style={font=\scriptsize,text=Muted,text width=2.9cm,align=center}]
 \node[stp] (s1) at (0,0) {\ico{\large\faRuler}\\[.05cm]\textbf{简单基线}\\\scriptsize 对照均值 + 平均效应};
 \node[stp] (s2) at (3.55,.4) {\ico{\large\faProjectDiagram}\\[.05cm]\textbf{GEARS}\\\scriptsize 2023 / 2024 · 基因关系图};
 \node[stp] (s3) at (7.1,.8) {\ico{\large\faLayerGroup}\\[.05cm]\textbf{基础模型微调}\\\scriptsize scGPT · scFoundation};
 \node[stp,fill=PaleTeal,draw=DarkBlue,line width=.9pt] (s4) at (10.65,1.2) {\ico{\large\faCubes}\\[.05cm]\textbf{State / Stack}\\\scriptsize 2026 · 细胞集合建模};
 \draw[flow] (s1)--(s2);\draw[flow] (s2)--(s3);\draw[flow] (s3)--(s4);
 \draw[-{Latex[length=2mm]},draw=Muted,line width=.6pt] (-1.55,2.25) -- node[above,font=\scriptsize,text=Muted]{模型复杂度、数据与算力需求递增} (12.2,2.25);
 \node[nt,below=.22cm of s1] {对照均值 + 扰动平均效应；双基因取加性。独立评测中常与深度模型持平。};
 \node[nt,below=.22cm of s2] {基因关系图上的图神经网络；面向同一背景内的新扰动。};
 \node[nt,below=.22cm of s3] {先预训练细胞表征再微调；微调能力不等于零样本泛化。};
 \node[nt,below=.22cm of s4] {在细胞集合上建模状态转移；Stack 用上下文学习；接近 VCC 2026 设置。};
\end{tikzpicture}
\remark{\ico{\faExclamationTriangle}\enspace 复杂度上升不保证增益：每一级都要与更简单的方法在同划分、同指标下比较。}
\source{[4--10,12,13]。所列模型均按原始论文或官方说明理解，尚未在本地运行。}
\end{frame}
